\section{Transfer from the reference volumes}\label{sec:uniform-transfer}

The renormalization estimates are obtained on two reference tori.
To use them for volume-first limits, we need bounds on products
of local observables in arbitrarily large cutoff tori. The
chessboard estimate reduces each such product to a folded
one-cell observable. A separate convexity estimate controls
the normalization by the partition functions of the two
reference volumes.

\subsection{Folded observables and the partition defect}

Use unit-lattice lengths in this subsection. Fix the charge $q>0$
and a real lattice mass. Choose coordinate cells of width
$w\in\tfrac12\Z_{>0}$ and offsets $o_\mu\in\tfrac12\Z$:
\[
 Q_c^\circ=\{y:o_\mu+wc_\mu<y_\mu<o_\mu+w(c_\mu+1),\
                           \mu=1,2,3\}.
\]
The walls are admissible coordinate mirrors. For even positive
integers $s_\mu$ with $ws_\mu\ge4$, let $C(ws)$ be the rectangular
torus with these side lengths and antiperiodic fermions.
It has $s_1s_2s_3$ cells. The wall-reflection group acts simply
transitively on the covering cells. Write $O^{[c]}$ for the
reflected copy of $O\in\mathcal A(Q_0^\circ)$ in $Q_c^\circ$,
and
\[
 O^\sharp_{ws}=\prod_{c\bmod s}O^{[c]},\qquad
 \mathcal Z_s(O)=Z_{C(ws)}[O^\sharp_{ws}].
\]
The copies use the antilinear reflection action, including
conjugation of scalar coefficients.

These arrays satisfy the rectangular chessboard estimate
\cite[(69)]{ChamberLattice}. The cell theorem is stated there for integer width.
For the coordinate cells used here, its fold proof also applies
at half-integer width: all walls remain admissible, paired
reflections translate by the integer $2w$, and the even cell
counts give integer torus sides at least four.
For a support box $B=y+[-r,r]^3\subset Q_c^\circ$ one has
$2r<w$. A nonzero period has sup norm at least $2w>4r$
and at least four. Thus its enlarged boxes $B+[-1,1]^3$
are disjoint, by the same embedding estimate used in the
proof of \cite[Proposition~4.3]{ChamberLattice}.
The products therefore belong to the required slab classes.

Hold two entries of $s$ fixed and vary $s_\mu$.
The product of the copies in the layer $c_\mu=0$ gives
a Hilbert--Schmidt slab operator $Y$ independent of $s_\mu$.
Its reflected layer pairs fill the entire folded array.
The slab-fold identity \cite[Proposition~4.2, (66)]{ChamberLattice} gives
\begin{equation}\label{eq:transfer-fold-trace}
 \mathcal Z_s(O)=\operatorname{Tr}(Y^*Y)^{s_\mu/2}.
\end{equation}
In particular $\mathcal Z_s(O)\ge0$, and
$s_\mu\mapsto\mathcal Z_s(O)^{1/s_\mu}$ is nonincreasing.
If $Y\ne0$, the logarithm of this folded partition function
is convex in $s_\mu$. The first assertion is the monotonicity
of Schatten norms, and the second follows from H\"older's
inequality for the trace of positive powers.
The normalization of $Y$ here is that of the slab operators
in \cite[\S4.3]{ChamberLattice}.

Choose an even reference cell count $s_0\ge2$ with
$n_r=ws_0\ge4$, and an integer $j\ge2$. Put
\begin{equation}\label{eq:transfer-reference-data}
 N_r=s_0^3,\qquad
 D_j=j^3\log Z_{C(n_r)}[1]-\log Z_{C(jn_r)}[1],
 \qquad \delta_j=\frac{D_j}{(j-1)^3N_r}.
\end{equation}
For an observable in the base cell its normalized reference
seminorm is
\begin{equation}\label{eq:transfer-reference-seminorm}
 \|O\|_*=
 \left(\frac{\mathcal Z_{(s_0,s_0,s_0)}(O)}
                    {Z_{C(n_r)}[1]}\right)^{1/N_r}.
\end{equation}
The partition function is positive, and the numerator is
nonnegative, so this is well-defined also when it vanishes.

\begin{proposition}\label{prop:transfer-general}
One has $D_j\ge0$. Let $s$ be even, with either $s=s_0$
or $s\ge js_0$, and put $r=s/s_0$, $n=ws$.
For every set $\mathcal C$ of distinct cells of $C(n)$ and
$O_c\in\mathcal A(Q_0^\circ)$,
\begin{equation}\label{eq:transfer-general}
 \left|\E_{C(n)}\prod_{c\in\mathcal C}O_c^{[c]}\right|
                    \le\prod_{c\in\mathcal C}\exp(\delta_j)\|O_c\|_*.
\end{equation}
The right side uses only the two reference partition functions
and the one-cell observables folded on $C(n_r)$.
\end{proposition}

\begin{proof}
Successively decreasing the three side counts in
\eqref{eq:transfer-fold-trace} gives, for $s\ge s_0$,
\begin{equation}\label{eq:transfer-fold-upper}
 \mathcal Z_{(s,s,s)}(O)
            \le\mathcal Z_{(s_0,s_0,s_0)}(O)^{r^3}.
\end{equation}
Taking $O=1$ and $r=j$ proves $D_j\ge0$.

We next obtain the lower bound for the normalizing partition
function. Define
$p(r_1,r_2,r_3)=\log Z_{C(r_1n_r,r_2n_r,r_3n_r)}[1]$
on ratios with integer side lengths at least four.
The pressure trace \cite[(75)]{ChamberLattice} shows that
$p$ is convex in each argument and that $p/r_\mu$ is
nonincreasing in that argument. For $r\ge j$ put
$a=(r-1)/(j-1)$. Convexity gives
$p(r,\cdot,\cdot)\ge a p(j,\cdot,\cdot)
-(a-1)p(1,\cdot,\cdot)$.
Apply this in each argument, retaining the positive coefficient
$a$ at the next step. Bound the three subtracted terms by the
side monotonicity. The result is
\begin{align*}
 p(r,r,r)\ge{}&
 a^3p(j,j,j)\\
 &-(a-1)(a^2j^2+ajr+r^2)p(1,1,1)\\
 ={}&r^3p(1,1,1)-a^3D_j.
\end{align*}
The last equality uses $aj-(a-1)=r$ three times.
The negative coefficients are why upper bounds are used on
the three subtracted terms; no sign condition on $p$ is needed.
Thus, including the equality case $r=1$,
\begin{equation}\label{eq:transfer-partition-lower}
 \log Z_{C(n)}[1]\ge
 r^3\left(\log Z_{C(n_r)}[1]-\frac{D_j}{(j-1)^3}\right).
\end{equation}

Fill the cells outside $\mathcal C$ with the constant integrand.
The normalized rectangular chessboard estimate then gives
\[
 \left|\E_{C(n)}\prod_{c\in\mathcal C}O_c^{[c]}\right|
 \le\prod_{c\in\mathcal C}
       \left(\frac{\mathcal Z_{(s,s,s)}(O_c)}
                         {Z_{C(n)}[1]}\right)^{1/(r^3N_r)}.
\]
Use \eqref{eq:transfer-fold-upper} in the numerator and
\eqref{eq:transfer-partition-lower} in the denominator.
Each factor is at most $\exp(\delta_j)\|O_c\|_*$, proving
\eqref{eq:transfer-general}.
\end{proof}

